% Generated draft: requires mathematical and source review.
\documentclass[11pt]{article}
\usepackage[margin=25mm]{geometry}
\usepackage{fontspec}
\setmainfont{Latin Modern Roman}
\usepackage{amsmath,amssymb,mathtools,unicode-math}
\setmathfont{Latin Modern Math}
\usepackage{enumitem,needspace,xurl,hyperref}
\setlength{\parindent}{0pt}
\setlength{\parskip}{5pt}
\setlength{\emergencystretch}{3em}
\newcommand{\N}{\mathbb{N}}
\newcommand{\Z}{\mathbb{Z}}
\newcommand{\Q}{\mathbb{Q}}
\newcommand{\R}{\mathbb{R}}
\newcommand{\C}{\mathbb{C}}
\newcommand{\F}{\mathbb{F}}
\newcommand{\A}{\mathbb{A}}
\newcommand{\PP}{\mathbb{P}}
\newcommand{\E}{\mathbb{E}}
\newcommand{\eps}{\varepsilon}
\newcommand{\dd}{\,\mathrm{d}}
\begin{document}
\section*{UnsolvedMath ID: 800936}
\textbf{Source problem code:} catalogue entry
\subsection{Definitions and mathematical statement}
Let \(X\) be a two-dimensional simplicial complex equipped with a metric \(d\) obtained by gluing hyperbolic triangles of curvature \(-1\) along geodesic edges. Assume that the metric is locally \(\operatorname{CAT}(-1)\) at every vertex and that every edge belongs to at least two two-dimensional simplices. Here local \(\operatorname{CAT}(-1)\) means that each point under consideration has a neighborhood satisfying the \(\operatorname{CAT}(-1)\) triangle-comparison inequality. The supplied question describes these complexes as locally \(\operatorname{CAT}(-1)\).

For a metric space \((Z,\rho)\), let \(\mathcal{C}(Z)\) denote its set of free homotopy classes of loops. One possible convention for its marked length spectrum is the function
\[
L_{\rho}\colon\mathcal{C}(Z)\longrightarrow[0,+\infty],\qquad
L_{\rho}(c)=\inf\{\operatorname{Length}_{\rho}(\gamma):\gamma\text{ is a rectifiable closed loop representing }c\},
\]
with \(\inf\varnothing=+\infty\). This convention does not assume that the infimum is attained by a closed geodesic. The source does not specify this convention or hypotheses guaranteeing such attainment.

\textbf{Question 16, attributed in the supplied record to Benjamin Schmidt.} Does the marked length spectrum determine \(d\) in either of the following comparison classes?
\begin{enumerate}
\item All locally \(\operatorname{CAT}(-1)\) metrics \(d'\) on the same space \(X\).
\item All metric two-dimensional simplicial complexes \((Y,d')\) homeomorphic to \(X\), whose simplices are hyperbolic with geodesic edges, whose vertex neighborhoods satisfy the local \(\operatorname{CAT}(-1)\) condition, and whose edges each belong to at least two two-dimensional simplices.
\end{enumerate}
Using the preceding spectrum convention, the spectral hypothesis in the first class is
\[
L_d(c)=L_{d'}(c)\qquad\text{for every }c\in\mathcal{C}(X).
\]
In the second class, a homeomorphism \(h\colon X\to Y\) supplies a marking, and the spectral hypothesis is
\[
L_d(c)=L_{d'}(h_*c)\qquad\text{for every }c\in\mathcal{C}(X),
\]
where \(h_*c\) is the free homotopy class obtained by composing a representative with \(h\).

A marking-preserving rigidity formulation asks whether these equalities imply the existence of an isometry \(F\colon(X,d)\to(Y,d')\) homotopic to \(h\); in the first class take \(Y=X\) and \(h=\operatorname{id}_X\). The supplied wording does not specify whether this marking-preserving conclusion, preservation only of the induced free homotopy classes, or merely existence of an isometry is intended. It also supplies no compactness, finiteness, local finiteness, connectedness, or completeness hypotheses. Consequently, the displayed formulation records a conventional interpretation and the unresolved qualifications, rather than an unqualified precise rigidity conjecture.
\subsection{Short English statement}
Benjamin Schmidt's question asks whether the lengths associated with all free homotopy classes of closed loops determine a locally \(\operatorname{CAT}(-1)\) hyperbolic triangle complex, either among metrics on the same space or among homeomorphic complexes of the stated type. The source leaves the exact rigidity conclusion and several global hypotheses unspecified.
\subsection{Sources}
\begin{itemize}
\item Local UnsolvedMath registry, permanent ID \(800936\), source problem code \texttt{catalogue entry}; supplied complete source record, upstream revision \texttt{d371ee917ac9647fbb06cd6fc529deec6f3dda82}.
\item Burns and Matveev, \emph{Open Problems and Questions about Geodesics}, Question 16, listed in the supplied record as a 2019 source: \url{https://arxiv.org/abs/1308.5417}. The record identifies the extraction method as \texttt{arxiv\_latex} and gives an access date of 2026-10-16. These are supplied provenance details; the source was not independently accessed or verified for this draft.
\item The registry lists Burns and Matveev as proposers and gives Benjamin Schmidt as the question's attribution/context. Its status label is \texttt{open}; this is a supplied classification, not an independently verified assessment of the literature.
\end{itemize}
\end{document}
